\documentclass[11pt,a4paper]{article}

% ============================================================
% Standard English mathematical typesetting preamble
% ============================================================
\usepackage[utf8]{inputenc}
\usepackage[T1]{fontenc}
\usepackage[english]{babel}
\usepackage{amsmath,amssymb,amsthm}
\usepackage{geometry}
\geometry{margin=1in}
\usepackage{graphicx}
\usepackage{booktabs}
\usepackage{array}
\usepackage{fancyhdr}
\usepackage{titlesec}
\usepackage{enumitem}
\usepackage{mathtools}
\usepackage{bm}

% Theorem-like environments
\theoremstyle{plain}
\newtheorem{theorem}{Theorem}
\newtheorem{lemma}{Lemma}
\theoremstyle{definition}
\newtheorem{remark}{Remark}

% Page style
\pagestyle{fancy}
\fancyhf{}
\fancyhead[C]{\small Cayley --- Collected Mathematical Papers, Vol.\ VIII}
\fancyfoot[C]{\thepage}
\renewcommand{\headrulewidth}{0.4pt}

% Section formatting to match original numbered articles
\titleformat{\section}{\large\scshape\centering}{}{0em}{}
\titleformat{\subsection}{\normalsize\bfseries}{}{0em}{}

% Custom commands for Cayley's notation
\newcommand{\dd}{\,d}
\DeclareMathOperator{\Jac}{Jac}

\begin{document}

% ============================================================
% PAPER 538
% ============================================================
\section{[538]}
\subsection{Extract from a Letter to Mr C.~W.~Merrifield.}

\begin{center}
\textit{[From the Messenger of Mathematics, vol.~\textsc{i}. (1872), pp.~87, 88.]}
\end{center}

\textsc{The} general integral of the equations
\[
\frac{p}{x} = \frac{dt}{a} = \frac{dq}{y}.
\]
\[
\left[ \text{where } a, \, \beta, \, \gamma, \, \delta = \frac{dq}{dx}, \, \frac{dq}{dy}, \, \frac{dq}{ds}, \, \frac{dq}{dt}, \right] \text{ can, I think, be found, viz.\ } \frac{p}{x} = \frac{dt}{a} = \frac{dq}{y}
\]
$s =$ function $x$, and $\frac{dt}{ds} = \frac{a}{y}$ gives $s =$ function $t$. But $s =$ function $x$, is integrated as the equation of a developable surface ($p$ instead of $s$), we must have
\[
p = ax + by + p',
\]
\[
0 = ax + y + p';
\]
similarly, $s =$ function $t$, gives
\[
\left( a' \frac{dx}{ds}, \quad y' \frac{dq}{dt} \right);
\]
$a$ and $p$ functions of $t$, and
\[
\left( a' = \frac{dx}{ds}, \quad y' = \frac{dq}{dt} \right);
\]

Observe that the constants here have been so taken, that $\frac{dq}{ds} = a$, $\frac{dq}{dt} = b$; in order that $\delta$ may, in the two pairs of equations, mean the same function of $(x, y)$, we must have
\[
s' = \frac{b}{a} - \frac{p}{x}.
\]

\clearpage

\bigskip
\textit{p.~518 [538]}
\smallskip

that is,
\[
b = \int \frac{dx}{y}, \quad t' = \int \frac{dq}{x};
\]
or, writing $a = dt$, $p = q$, we have
\[
p = abh + pb + q,
\]
\[
q = bs + p \int \frac{dx}{y} - p \int \frac{a}{y} \, dx,
\]
where
\[
abh + y + qh = 0.
\]

The last equation gives $b$ as a function of $(x, y)$, and the values of $p$, $q$ are then such that $dx + pdx + qdy$ is a complete differential, so that we obtain $z$ by the integration of that equation.

A simple example is
\[
p = \tfrac{1}{2} bs - bq, \quad q = -bs + g \log b, \quad bs - y = 0,
\]
that is,
\[
p = -\tfrac{1}{4} \frac{y^2}{x}, \quad q = -y + y \log \frac{y}{x},
\]
whence
\[
z = \tfrac{1}{4} y^2 \log \frac{y}{x} - \tfrac{1}{4} y^2;
\]
we have
\[
r = \tfrac{1}{2} \frac{y^2}{x^2}, \quad s = -\frac{y}{x}, \quad t = \log \frac{y}{x},
\]
\[
a = -\frac{y^2}{2x^2}, \quad \beta = -\frac{y}{x}, \quad \gamma = -\frac{1}{y}, \quad \delta = \frac{1}{x};
\]
or
\[
\frac{a}{\beta} = \frac{p}{x} = \frac{b}{y} \left( = a - \frac{p}{x} \right);
\]
as it should be.

Cambridge, 28 July, 1871.

\clearpage

% ============================================================
% PAPER 539
% ============================================================
\section{[539]}
\subsection{Further Note on Lagrange's Demonstration of Taylor's Theorem.}

\begin{center}
\textit{[From the Messenger of Mathematics, vol.~\textsc{i}. (1872), pp.~105, 106.]}
\end{center}

This refers to a paper, Wilkinson, ``Note on Taylor's Theorem,'' \textit{Messenger of Mathematics}, same volume, pp.~96, 97, discussing the ``Note on Lagrange's Demonstration of Taylor's Theorem,'' [248, ante, pp.~493---495].

\clearpage

% ============================================================
% PAPER 540
% ============================================================
\section{[540]}
\subsection{On a Property of the Torse Circumscribed about Two Quadric Surfaces.}

\begin{center}
\textit{[From the Messenger of Mathematics, vol.~\textsc{i}. (1872), pp.~111, 112.]}
\end{center}

\textsc{The} property mentioned by Mr Townsend in his paper(*) in the August No., ``On a Property in the Theory of Confocal Quadrics,'' may be demonstrated in a form which, it appears to me, better exhibits the foundation and significance of the theorem.

Starting with two given quadric surfaces, the torse circumscribed about these touches each of a singly infinite series of quadric surfaces, any two of which may be used (instead of the two given surfaces) to determine the torse; in the series are included four conics, one of them in each of the planes of the self-conjugate tetrahedron of the two given surfaces; and if we attend to only two of these conics, the two conics are in fact any two conics whatever, and the torse is the circumscribed torse of the two conics; or, what is the same thing, it is the envelope of the common tangent-planes of the two conics.

Consider new two conics $U$, $U'$, the planes of which intersect in a line $I$; and let $I$ meet $U$ in the points $L$, $M$, and meet $U'$ in the points $L'$, $M'$: take $L$ the pole of $I$ in regard to the conic $U$, and $L'$ the pole of $I$ in regard to the conic $U'$.

Take $T$ any point on $I$, and draw $TP$ touching $U$ in $P$, and $TP'$ touching $U'$ in $P'$: the points $P$, $P'$ may be considered as corresponding points on the two conics.

Join $AP$ and produce it to meet the line $I$ in $G$; the line $APG$ is in fact the polar of $T$ in regard to the conic $U$ (for $T$ being a point on $I$, the polar of $T$ passes through $A$; and this polar also passes through $P$); that is, the points $T$, $G$, and $L$, $M$ are harmonics on the line $I$; whence also, in the plane of the conic $U'$,

\smallskip
\hrule height 0.4pt
\smallskip
\textsuperscript{(*)} \textit{Messenger of Mathematics, same volume, pp.~48, 49.}

\clearpage

\bigskip
\textit{p.~520 [540]}
\smallskip

the lines $PG$, $PT$ and $PL$, $PM$ are harmonic lines through the point $P$. It thus appears that in the particular case where the points $L$, $M$ are the foci of the conic $U'$, the line $PG$ is the normal at the point $P$; and we may say in general that $PG$ is the quasi-normal at the point $P$ of the conic $U'$.

\begin{center}
\textit{[Figure: Conic with quasi-normal construction --- diagram omitted]}
\end{center}

Consider now the torse circumscribed about the conics $U$, $U'$; the plane $PTP'$ will represent any plane, and the line $PP'$ any line of this torse; projecting on the plane of $U'$ with the point $A$ as centre of projection, the projection of $PP$ is the line $PG$; which, as just seen, is the quasi-normal of the conic $U'$ at the point $P'$.

The projection of the cuspidal curve is the envelope of line $PG$, which is the projection of the generating line $PP'$ of the torse---viz.\ this envelope is the quasi-evolute of the conic $U'$; which is the theorem in question.

\clearpage

% ============================================================
% PAPER 541
% ============================================================
\section{[541]}
\subsection{On the Reciprocal of a Certain Equation of a Conic.}

\begin{center}
\textit{[From the Messenger of Mathematics, vol.~\textsc{i}. (1872), pp.~120, 121.]}
\end{center}

\textsc{The} following formula is useful in various problems relating to conics: the reciprocal equation of the conic
\[
\lambda (px + qy + cz)(x + by + dz) + \mu (a''x + b''y + c''z)(a''x + b''y + c''z) = 0
\]
may be written indifferent in either of the forms
\[
|\lambda \, \xi\,, \eta\,, \zeta \, = a \, \xi\,, \eta\,, \zeta \, , \, b \, \xi\,, \eta\,, \zeta \, \, |+ 4 \lambda \mu \, \xi\,, \eta\,, \zeta \, , \, a' \, \xi\,, \eta\,, \zeta \, , \, b' \, \xi\,, \eta\,, \zeta \, \, |+ \, |a' \, \xi\,, \eta\,, \zeta \, , \, b' \, \xi\,, \eta\,, \zeta \, , \, c' \, \xi\,, \eta\,, \zeta \, \, | = 0,
\]
and
\[
|\lambda \, \xi\,, \eta\,, \zeta \, = a \, \xi\,, \eta\,, \zeta \, , \, b \, \xi\,, \eta\,, \zeta \, \, |+ 4 \lambda \mu \, \xi\,, \eta\,, \zeta \, , \, a \, \xi\,, \eta\,, \zeta \, , \, b'' \, \xi\,, \eta\,, \zeta \, \, |+ \, |a'' \, \xi\,, \eta\,, \zeta \, , \, b'' \, \xi\,, \eta\,, \zeta \, , \, c'' \, \xi\,, \eta\,, \zeta \, \, | = 0.
\]

In fact, in the reciprocal equation, noting the coefficient of $\beta$, it is
\[
\{\lambda (bc' - b'c) + \mu (b''c' - b''c')\}^2 = (Dab' - 2ab' b'') (Dac' - 2bc'c'),
\]
viz.\ this is
\[
\lambda^2 (bc' - b'c)^2 + \mu^2 (b''c' - b''c')^2 + 2\lambda\mu \left\{ \begin{array}{l} 2bb'c''c'' + 2b''b'cc' \\ -(bc + b'c')(b''c'' + b''c'') \end{array} \right\}.
\]

\clearpage

\bigskip
\textit{p.~523 [541]}
\smallskip

or, as it may be written,
\[
[\lambda (bc' - b'c) \pm \mu (b''c' - b''c')]^2 = D\lambda\mu \left\{ \begin{array}{l} 2bb'c''c'' + 2b''b'cc' \\ -(bc + b'c')(b''c'' + b''c'') \end{array} \right\}.
\]

Taking the upper signs, this is
\[
[\lambda (bc' - b'c) + \mu (b''c' - b''c')]^2 + 4\lambda\mu \left\{ \begin{array}{l} bb'c''c'' + b''b'cc', \\ -bc'b''c'' - b'cb''c'' \end{array} \right\},
\]
viz.\ the term in $\lambda\mu$ is
\[
+ 4\lambda\mu (bc'' - b''c)(b'c' - b'c''),
\]

Taking the lower signs, it is
\[
[\lambda (bc' - b'c) - \mu (b''c' - b''c')]^2 + 4\lambda\mu \left\{ \begin{array}{l} bb'c''c'' + b''b'cc', \\ -bc'b''c'' - b'cb''c'' \end{array} \right\},
\]
viz.\ the term in $\lambda\mu$ is
\[
+ 4\lambda\mu (bc' - b''c)(b'c'' - b''c').
\]

And it is hence easy to infer the forms of the other coefficients, and to arrive at the foregoing result.

\clearpage

\end{document}
